\section{Einrichtung des Testsystems}\label{appendix1:Einrichtung_des_Testsystems}
Zur Einrichtung des Testsystems wird zuallererst ein Linux-Container eingerichtet, welches ein einfacheres und schnelleres Deployen auf unterschiedlichen Linux-Distributionen ermöglicht. Zum Hosten des Linux-Containers wird in dieser Installationsanleitung Ubuntu 20.04 Focal Fossa verwendet, welches eine LTS-Version von Ubuntu darstellt. Nach der Einrichten des Linux-Containers können die benötigten Pakete in dieser neu erstellten Umgebung installiert und daraufhin entsprechend konfiguriert werden.

\subsection{Erstellen des Linux-Containers}
Für die Installation der Linux-Container-Tools werden Administratorrechte auf dem Zielsystem benötigt. Die benötigten Installationsschritte zum Einrichten des Linux-Containers sind wie folgt \cite{lxd_blog_nodate}.

\begin{enumerate}[label*=\arabic*.]
    \item LXD 4.0 installieren:\
        \begin{lstlisting}[language=bash]
sudo apt install lxd
        \end{lstlisting}
    bzw.
        \begin{lstlisting}[language=bash]
sudo snap install lxd
        \end{lstlisting}

    \item LXD initialisieren:\
        \begin{lstlisting}[language=bash]
lxd init
        \end{lstlisting}

    \item Den derzeitigen User zur LXD-Gruppe hinzufügen:\
        \begin{lstlisting}[language=bash]
sudo adduser <user name> lxd
        \end{lstlisting}
        Logout \& Login durchführen oder mittels des folgenden Befehls den User im Terminal neu anmelden.
        \begin{lstlisting}[language=bash]
su - <user name>
        \end{lstlisting}

    \item Erstellen des Linux-Containers:\\
        Der Container wird hiermit im Default-Container-Pool erstellt.
        \begin{lstlisting}[language=bash]
lxc init ubuntu:18.04 <container name>
        \end{lstlisting}
        Mit der Option \quoteMark{-s} kann ein anderer Container-Pool ausgewählt werden.
        \begin{lstlisting}[language=bash]
lxc init ubuntu:18.04 <container name> -s <container pool name>
        \end{lstlisting}

    \item Die gespeicherte Disk kann mit dem folgenden Befehl angeschaut werden:\
        \begin{lstlisting}[language=bash]
lxc storage list
        \end{lstlisting}

    \item Dem Linux-Container erlauben den X11-Traffic weiterzuleiten:
        \begin{enumerate}[label*=\arabic*.]
            \item Die Datei \quoteMark{gui-profile.yaml} mit dem folgenden Inhalt erstellen:
                \begin{lstlisting}[language=yaml]
config:
  environment.DISPLAY: :0
  raw.idmap: both 1000 1000
  user.user-data: |
    #cloud-config
    runcmd:
      - 'sed -i "s/; enable-shm =
 yes/enable-shm = no/g" /etc/pulse/client.conf'
      - 'echo export PULSE_SERVER=unix:/tmp/
.pulse-native | tee --append /home/ubuntu/.profile'
    packages:
      - x11-apps
      - mesa-utils
      - pulseaudio
description: GUI LXD profile
devices:
  PASocket:
    path: /tmp/.pulse-native
    source: /run/user/1000/pulse/native
    type: disk
  X0:
    path: /tmp/.X11-unix/X0
    source: /tmp/.X11-unix/X0
    type: disk
  mygpu:
    type: gpu
name: gui
used_by:
                \end{lstlisting}

                \item Bei der Erstellung der Datei \quoteMark{gui-profile.yaml} ist zu beachten, dass die IDs im Eintrag \quoteMark{raw.idmap: both 1000 1000} folgenden Bedeutung haben:
                    \begin{lstlisting}[language=yaml]
raw.idmap: both <host user id> <user id in container>
                    \end{lstlisting}

                    \begin{tabular}{|>{\raggedright\arraybackslash}p{0.18\textwidth}|>{\raggedright\arraybackslash}p{0.65\textwidth}|}
                        \hline
                        <host user id>: & Die User-ID, welche zum Starten des Linux-Containers verwendet wird.\\
                        \hline
                        <user id in container>: & Die User-ID des Users innerhalb des Linux-Containers. Normalerweise ist dies die User-ID 1000, da der erste User, welcher im Linux-Container eingerichtet wurde, standardmäßig die User-ID 1000 erhält.\\
                        \hline
                    \end{tabular}\\

                \item Das Profil kann daraufhin über die Datei \quoteMark{gui-profile.yaml} geladen werden:
                    \begin{lstlisting}[language=bash]
lxc profile create gui
cat gui-profile.yaml | lxc profile edit gui
                    \end{lstlisting}
        \end{enumerate}

    \item Den Linux-Container starten und eine Shell im Linux-Container öffnen:
        \begin{lstlisting}[language=bash]
lxc start <container name>
lxc exec <container name> -- sudo --user ubuntu --login
        \end{lstlisting}

    \item Um Rendering-Probleme zu vermeiden sollten die neuesten Treiber im Linux-Container installiert werden:
        \begin{lstlisting}[language=bash]
sudo add-apt-repository ppa:oibaf/graphics-drivers
sudo apt upgrade -y
        \end{lstlisting}

    \item Um die durchgeleitete Grafikkarte zu testen muss der Linux-Container neu gestartet werden! Anschließend müssen die folgenden Befehle eingegeben werden:
        \begin{lstlisting}[language=bash]
lspci | grep VGA
        \end{lstlisting}
        bzw.
        \begin{lstlisting}[language=bash]
sudo apt-get install glmark2
glmark2
        \end{lstlisting}

        \begin{redbox}{Wichtig:}
            Anstatt \quoteMark{llvmpipe} sollte die jeweilige verbaute Grafikkarte angezeigt werden.
        \end{redbox}

    \item Für die Soundausgabe sollte PulseAudio getestet werden. Es sollte beim Ausführen des folgenden Befehls ein Rauschen zu hören sein:
        \begin{lstlisting}[language=bash]
pacat < /dev/urandom
        \end{lstlisting}

    \item Von Zeit zu Zeit sollte ein Backup vom Linux-Container erstellt werden
        \begin{lstlisting}[language=bash]
lxc export ros-melodic ros-melodic.tar.gz
        \end{lstlisting}

        bzw. zum Erstellen eines Snapshots

        \begin{lstlisting}[language=bash]
lxc snapshot ros-melodic snap01
lxc publish ros-melodic/snap01 --alias <new-alias>
lxc image export <new-alias> ros-melodic.tar.gz
        \end{lstlisting}

    \item Das Backup kann mittels des folgenden Befehls zurückgespielt werden:
        \begin{lstlisting}[language=bash]
lxc import <new-alias> ros-melodic.tar.gz
        \end{lstlisting}

        bzw.

        \begin{lstlisting}[language=bash]
lxc image import ros-melodic.tar.gz --alias <new-alias>
        \end{lstlisting}

        \begin{redbox}{Wichtig:}
            Beim Import des Linux-Containers muss sichergestellt werden, dass zuvor alle mit dem Linux-Container verknüpften Linux-Container-Profile vorhanden sind, da es ansonsten zu Import-Fehlern kommt.
        \end{redbox}
\end{enumerate}


\subsection{Installation \& Konfiguration der Pakete}
Mit der zuvor installierten \quoteMark{Ubuntu 20.04 Focal Fossa}-Umgebung können nun die für diese Arbeit benötigten Pakete innerhalb dieser Umgebung installiert und konfiguriert werden.

\begin{enumerate}[label*=\arabic*.]
    \item Installation von ROS 1:\\
    Zur Nutzung von ROS 1 muss ROS Melodic installiert werden. Die Installationsanleitung kann der offiziellen ROS-Seite (\href{http://wiki.ros.org/melodic/Installation/Ubuntu}{http://wiki.ros.org/melodic/Installation/Ubuntu}) entnommen werden. Bei der Installation des ROS-Pakets muss nur das Basis-ROS-Paket installiert werden.

        \begin{lstlisting}[language=bash]
sudo apt install ros-melodic-ros-base
        \end{lstlisting}

    \item Installation des WidowX 200:
        \begin{enumerate}[label*=\arabic*.]
            \item Hierzu müssen zuvor die benötigten Abhängigkeiten installiert werden:

                \begin{lstlisting}[language=bash]
sudo apt install ros-melodic-robot-state-publisher ros-\
melodic-rqt-plot ros-melodic-joy
                \end{lstlisting}

            \item Anschließend müssen die auf der Webseite \href{https://github.com/Interbotix/interbotix_ros_arms#quickstart}{https://github.com\\/Interbotix/interbotix\_ros\_arms\#quickstart} beschriebenen Schritte durchgeführt werden.

            \item Danach müssen die folgenden Einstellungen auf dem Host-System durchgeführt werden damit der unprivilegierte Linux-Container Zugriff auf den Dynamixel U2D2 erhält: \textcolor{red}{\newline\newline TODO: eigenes repository oder CD\newline\newline}

                \begin{lstlisting}[language=bash]
wget https://raw.githubusercontent.com/Interbotix/interbotix_\
ros_arms/master/interbotix_sdk/10-interbotix-udev.rules
sudo cp 10-interbotix-udev.rules /etc/udev/rules.d/
sudo udevadm control --reload-rules && sudo udevadm trigger
echo "# Interbotix WidowX 200
sudo chown ubuntu /dev/ttyDXL" >> ~/.bashrc
                \end{lstlisting}

            \item Außerhalb des Containers müssen folgende Einstellungen erfolgen, damit der unprivilegierte Container Zugriff auf den Dynamixel U2D2 erhält:
                \begin{lstlisting}[language=bash]
lxc config device add <container name> dynamixel_u2d2_ttydxl unix-char \
path=/dev/ttyDXL required=false
                \end{lstlisting}
        \end{enumerate}

    \item Testen des WidowX 200:
        \begin{redbox}{Wichtig:}
            \begin{compactitem}
                \item Genügend Abstand einhalten
                \item Den Roboter festschrauben
                \item Beim Beenden der Softwaretools den Roboter festhalten, da dieser ansonsten in sich zusammenfällt.
            \end{compactitem}
        \end{redbox}

        Zum Testen der grundlegenden Funktionalität des WidowX 200 können ein paar spezifische ROS-Befehle an diesen gesendet werden.

        \begin{lstlisting}[language=bash]
roslaunch interbotix_descriptions description.launch robot_name:=\
wx200 jnt_pub_gui:=true

roslaunch interbotix_sdk arm_run.launch robot_name:=wx200

rosservice call /wx200/torque_joints_on

roslaunch interbotix_gazebo gazebo.launch robot_name:=wx200
rosservice call /gazebo/unpause_physics
        \end{lstlisting}

    \item Wichtige Informationen für ROS bekanntmachen:
        \begin{lstlisting}[language=bash]
echo "# ROS
source /opt/ros/melodic/setup.bash
source ~/interbotix_ws/devel/setup.bash" >> ~/.bashrc
source ~/.bashrc
        \end{lstlisting}

    \item Azure Kinect installieren:
        \begin{enumerate}[label*=\arabic*.]
            \item Außerhalb des Linux-Containers müssen die folgenden Einstellungen erfolgen um die Azure Kinect für den derzeitgen User verfügbar zu machen \cite{microsoftazure-kinect-sensor-sdk_installation_nodate}:

                \begin{lstlisting}[language=bash]
wget https://raw.githubusercontent.com/microsoft/Azure-Kinect-\
Sensor-SDK/develop/scripts/99-k4a.rules
sudo cp 99-k4a.rules /etc/udev/rules.d/
sudo udevadm control --reload-rules && sudo udevadm trigger
                \end{lstlisting}

            \item Außerhalb des Linux-Containers müssen die folgenden Einstellungen erfolgen, damit der unprivilegierte Container Zugriff auf die Azure Kinect erhält:
                \begin{lstlisting}[language=bash]
lxc config device add <container name> microsoft_generic_superspeed_usb_\
hub unix-char vendorid=045e productid=097a
lxc config device add <container name> microsoft_generic_usb_hub unix-char \
vendorid=045e productid=097b
lxc config device add <container name> azure_kinect_depth_camera unix-char \
vendorid=045e productid=097c
lxc config device add <container name> azure_kinect_4k_camera unix-char \
vendorid=045e productid=097d
lxc config device add <container name> azure_kinect_microphone_array \
unix-char vendorid=045e productid=097e
                \end{lstlisting}

            \item Den Linux-Container neustarten.
        \end{enumerate}

    \item Azure Kinect SDKs installieren:
        \begin{enumerate}[label*=\arabic*.]
            \item Microsoft-Paketquelle hinzufügen \cite{microsoftazure-kinect-sensor-sdk_installation_nodate}:
                \begin{lstlisting}[language=bash]
sudo apt install curl
curl https://packages.microsoft.com/keys/microsoft.asc | \
sudo apt-key add -
sudo apt-add-repository https://packages.microsoft.com/ubuntu\
/18.04/prod
                \end{lstlisting}

            \item Installieren der Azure Kinect SDKs:
                \begin{lstlisting}[language=bash]
sudo apt install k4a-tools=1.3.0
sudo apt install libk4a1.3-dev
                \end{lstlisting}

                \begin{redbox}{Wichtig:}
                    Die jeweilige Development-Version von libk4a installieren!

                    \begin{lstlisting}[language=bash]
sudo apt install libk4a<major>.<minor>-dev
                    \end{lstlisting}

                    z.B.

                    \begin{lstlisting}[language=bash]
sudo apt install libk4a1.3-dev
                    \end{lstlisting}
                \end{redbox}

            \item Um Paketkollisionen beim Updaten zu verhindern sollten das Paket \quoteMark{k4a-tools} vom Updateprozess ausgeschlossen werden:
                \begin{lstlisting}[language=bash]
sudo apt-mark hold k4a-tools
sudo apt-mark showhold
                \end{lstlisting}

            \item Nun kann das Paket \quoteMark{libk4abt1.0-dev} ohne Paketkollisionen installiert werden:
                \begin{lstlisting}[language=bash]
sudo apt install libk4abt1.0-dev
                \end{lstlisting}
        \end{enumerate}

    \item Azure Kinect Firmware aktualisieren:
        \begin{enumerate}[label*=\arabic*.]
            \item Überprüfen ob die Azure Kinect erkannt wird:
                \begin{lstlisting}[language=bash]
AzureKinectFirmwareTool -l
                \end{lstlisting}

            \item Die Ausgabe sollte in etwa wie folgt aussehen:
                \begin{lstlisting}[language=bash]
 == Azure Kinect DK Firmware Tool ==
Found 1 connected devices:
1: Device "000846393812"
                \end{lstlisting}

            \item Die jeweilige Firmware von \href{https://github.com/microsoft/Azure-Kinect-Sensor-SDK/blob/develop/docs/usage.md#msis}{https://github.com/microsoft/Azure-Kinect-Sensor-SDK\\/blob/develop/docs/usage.md\#msis} herunterladen und flashen
                \begin{lstlisting}[language=bash]
wget https://download.microsoft.com/download/1/9/8/198048e8-\
63f2-45c6-8f96-1fd541d1b4bc/AzureKinectDK_Fw_1.6.102075014.bin
                \end{lstlisting}

            \item Die heruntergeladene Firmware flashen:
                \begin{lstlisting}[language=bash]
AzureKinectFirmwareTool -u AzureKinectDK_Fw_1.6.102075014.bin
                \end{lstlisting}
        \end{enumerate}

    \item Azure Kinect testen:
        \begin{enumerate}[label*=\arabic*.]
            \item Device-Stream testen:
                \begin{lstlisting}[language=bash]
k4aviewer
                \end{lstlisting}

            \item Body Tracking testen: \textcolor{red}{Patch already in Master?}
                \begin{lstlisting}[language=bash]
sudo apt install mesa-utils
glxinfo | grep "OpenGL version"
                \end{lstlisting}

                \begin{redbox}{Wichtig:}
                    Das Body Tracking benötigt mindestens OpenGL 4.4 mit Compatibility Context.
                \end{redbox}

                Wenn eine Nvidia GPU verbaut ist, kann die GPU-Beschleunigung eingesetzt werden. Hierzu muss das CUDA-Paket installiert werden:
                \begin{lstlisting}[language=bash]
sudo apt install nvidia-cuda-toolkit
k4abt_simple_3d_viewer
                \end{lstlisting}

                Wenn keine Nvidia GPU verbaut ist muss ansonsten der CPU-Modus verwendet werden:
                \begin{lstlisting}[language=bash]
k4abt_simple_3d_viewer CPU
                \end{lstlisting}
        \end{enumerate}

    \item Zum Builden der Roboter-Gesten-Anwendung müssen die folgenden Pakete installiert werden:
        \begin{lstlisting}[language=bash]
sudo apt install cmake gdb gcc python python-catkin-tools
        \end{lstlisting}

    \item Source Code kompilieren:
        \begin{enumerate}[label*=\arabic*.]
            \item InterbotiX SDK kompilieren:
                \begin{lstlisting}[language=bash]
cd ~/interbotix_ws
catkin build
source ~/.bashrc
                \end{lstlisting}

            \item Die Roboter-Gesten-Anwendung kompilieren:
                \begin{lstlisting}[language=bash]
echo "# ROS
source ~/tir_ws/devel/setup.bash" >> ~/.bashrc
cd ~/tir_ws
catkin build
source ~/.bashrc
                \end{lstlisting}
        \end{enumerate}

    \item Entwicklungsumgebung installieren:
        \begin{enumerate}[label*=\arabic*.]
            \item Temporärer Ordner erstellen:
                \begin{lstlisting}[language=bash]
echo "# Temporary user folder
sudo mkdir -p /run/user/$(id --user)
sudo chown $(whoami) /run/user/$(id --user)" >> ~/.bashrc
                \end{lstlisting}

            \item VS Code installieren:
                \begin{lstlisting}[language=bash]
sudo snap install code --classic
                \end{lstlisting}

            \item Starten der Entwicklungsumgebung mit dem folgenden Befehl:
                \begin{lstlisting}[language=bash]
code
                \end{lstlisting}

            \item Mittels Strg + P in VS Code das \quoteMark{Quick Open}-Menü öffnen und jeweils eine der folgenden Befehle eingeben um die notwendigen Plugins zu installieren:
                \begin{lstlisting}[language=bash]
ext install ms-vscode.cpptools
ext install ms-vscode.cmake-tools
ext install ms-python.python
ext install ms-iot.vscode-ros
                \end{lstlisting}
        \end{enumerate}

        \item Wenn der Speicherplatz des Linux-Containers zu knapp wird kann der Speicherplatz des Linux-Containers wie folgt auf dem Host-System vergrößert werden:
            \begin{enumerate}[label*=\arabic*.]
                \item Image des Linux-Containers vergrößern:
                    \begin{lstlisting}[language=bash]
sudo truncate -s +5G /var/snap/lxd/common/lxd/disks/default.\
img
                    \end{lstlisting}

                \item Host-System neustarten

                \item Anschließend das Image des Linux-Containers einhängen und daraufhin die neue Größe des Images eingeben:
                    \begin{lstlisting}[language=bash]
sudo mount /var/snap/lxd/common/lxd/disks/default.img \
/<mount point>
sudo btrfs filesystem resize maxG /<mount point>
                    \end{lstlisting}
            \end{enumerate}
\end{enumerate}


\subsection{Starten der Roboter-Gesten-Anwendung}
Für die Roboter-Gesten-Anwendung wird als Scheduler SCHED\_DEADLINE eingesetzt, welcher für diese Arbeit feste Echtzeitfähigkeit im 50 ms Bereich garantieren kann.

\begin{figure}[htb]
	\centering
	\includegraphics[width=0.85\textwidth]{images/anhang/sched_deadline}
	\caption[Ablauf von SCHED\_DEADLINE]{Ablauf von SCHED\_DEADLINE \\Quelle: \cite{man_sched7_nodate}}
	\label{fig:sched_deadline}
\end{figure}
\FloatBarrier

Für den \quoteMark{SCHED\_DEADLINE}-Scheduler müssen hierfür die Parameter, welche in Abbildung \ref{fig:sched_deadline} ersichtlich sind, richtig eingestellt werden. Der \quoteMark{Runtime}-Parameter beträgt zur Einhaltung der 50 ms Frist hierbei 35000000 ns, der \quoteMark{Deadline}-Parameter beträgt 50000000 ns und der \quoteMark{Periode}-Parameter beträgt 50000000 ns. Der \quoteMark{Runtime}-Parameter steht hierbei für die maximal benötigte Ausführungszeit der Aufgabe, der \quoteMark{Deadline}-Parameter für wie lange genügend Zeit vorhanden ist um die Aufgabe auszuführen und der \quoteMark{Periode}-Parameter für wann die nächste Periode gestartet werden muss \cite{man_sched7_nodate}.\\

Die Roboter-Gesten-Anwendung kann mit dem Start-Script in Listing \ref{lst:start_script_for_app} auf dem Host-System im Linux-Container gestartet und dem \quoteMark{SCHED\_DEADLINE}-Scheduler zugewiesen werden.

\begin{lstlisting}[language=bash, caption={Script zum Starten der Roboter-Gesten-Anwendung}, label={lst:start_script_for_app}]
USE_ROS_COMMUNICATION=true
USER_NAME="ubuntu"
LXC_INSTANCE="<container name>"
ROS_APP_PARAMETER=""
ADDITIONAL_APP_PARAMETERS="--move-home-at-exit --move-home-at-error"
APP_MODE_PARAMETER="--teach-positions"  # --repeat-recorded-positions or
                                        # --teach-positions

lxc start $LXC_INSTANCE
sleep 3
USER_ID=$(lxc exec $LXC_INSTANCE -- id -u $USER_NAME)

if [ "$USE_ROS_COMMUNICATION" = true ] ; then
    ROS_APP_PARAMETER="--use-ros"
    # Start the InterbotiX node
    lxc exec $LXC_INSTANCE -- sudo -S -u $USER_NAME -i /bin/bash -i -c \
    "nohup roslaunch interbotix_sdk arm_run.launch robot_name:=wx200 &>\
    /dev/null &"
    sleep 3
else
    # Start roscore
    lxc exec $LXC_INSTANCE -- sudo -S -u $USER_NAME -i /bin/bash \
    -i -c "nohup roscore &>\
    /dev/null &"
    sleep 3
    # Initialize the parameter server
    lxc exec $LXC_INSTANCE -- sudo -S -u $USER_NAME -i /bin/bash \
    -i -c "nohup roslaunch interbotix_sdk arm_run.launch robot_name:=\
    wx200 &>/dev/null &"
    INTERBOTIX_PID=$(ps aux | grep -i "roslaunch interbotix_sdk" | \
    awk '{print $2}' | head -n -1)
    sleep 3
    # The InterbotiX node isn't required anymore
    kill $INTERBOTIX_PID
fi

# Start the TIR app
bash lxc exec $LXC_INSTANCE -- sudo -S -u $USER_NAME -i /bin/bash \
-i -c \
"rosrun teach_industrial_robots teach_industrial_robots_node \
$APP_MODE_PARAMETER $ROS_APP_PARAMETER $ADDITIONAL_APP_PARAMETERS" \
& LXC_TIR_APP_PID=$1

sleep 2
# Schedule the TIR app for real time usage
TIR_PID_APP=$(ps aux | grep -i "teach_industrial_robots/teach_\
industrial_robots_node" | \
awk '{print $2}' | head -n -1)
pkexec chrt -d --sched-runtime 1000000 --sched-deadline 5000000 \
--sched-period 5000000 -p  0 $TIR_PID_APP

if [ $(chrt -p $TIR_PID_APP | grep -c "SCHED_DEADLINE") -eq 0 ] ; \
then
    lxc stop $LXC_INSTANCE --force
    exit
fi

# Set the signal handlers
trap "lxc stop $LXC_INSTANCE --force" SIGINT SIGTERM SIGHUP

# Wait for the TIR app to complete
wait $LXC_TIR_APP_PID
\end{lstlisting}


% https://access.redhat.com/solutions/3742421

% lxc start script for app and set scheduler for process
% how to detect if lxc is running

%-----------------------------------------------

% LXC: https://serverfault.com/questions/630220/how-do-i-configure-lxc-to-allow-the-use-of-sched-rr-in-a-container
% Scheduler wechseln:
% cat /sys/block/sda/queue/scheduler
% https://www.thomas-krenn.com/de/wiki/Linux_I/O_Scheduler#Deadline


%--------------
%sudo apt install schedtool # http://manpages.ubuntu.com/manpages/xenial/man8/schedtool.8.html
%
%"A privileged user can change the priority policy of a process with the schedtool program[7]:ln 326, 373 or it is done by a program itself.[7]:ln 336 The priority class can be manipulated at the code level with a syscall like sched_setscheduler only available to root,[11] which schedtool uses.[12]"
%# https://en.wikipedia.org/wiki/Brain_Fuck_Scheduler
%sudo chrt -p $(pidof -s bash)


% "the default scheduler is CFS, the "Completely Fair Scheduler"
% http://manpages.ubuntu.com/manpages/focal/en/man7/sched.7.html
% http://manpages.ubuntu.com/manpages/eoan/en/man7/sched.7.html
% https://en.wikipedia.org/wiki/Completely_Fair_Scheduler
% https://www.kernel.org/doc/html/latest/scheduler/sched-design-CFS.html
%-------------

% https://www.researchgate.net/publication/331290349_The_real-time_linux_kernel_A_survey_on_Preempt_RT

% -----------


% sudo chrt -d --sched-runtime 1000000 --sched-deadline 5000000 --sched-period 5000000 -p  0 57802
% https://access.redhat.com/solutions/3742421
% http://manpages.ubuntu.com/manpages/cosmic/de/man1/chrt.1.html
% https://lwn.net/Articles/743740/
% # nanoseconds for the parameters

% TODO: Wie lange ist für uns Echtzeit? -> Messungen! und Begründung!

%-----------------------------------------------


% Dead Man's Switch: irgendetwas in der Hand halten?
% Wenn sich keine Person vor dem Tiefensesor befindet, dann werden auch keine Bewegungen ausgeführt.

% verwendeter Tiefensensor-Modus -> warum?

% Intel Real Sense Camera ZR300

%-------------------


% https://books.google.at/books?id=vVMoDwAAQBAJ&pg=PA229&lpg=PA229&dq=lxc+realtime&source=bl&ots=6aSRy33GKg&sig=ACfU3U13TeQVx0D1fZ5Z6DxLmKdeHTkXOg&hl=de&sa=X&ved=2ahUKEwj3mrWgyuvpAhUEGuwKHdcmDr8Q6AEwA3oECAoQAQ#v=onepage&q=lxc%20realtime&f=false

\section{Einrichtung der Simulationsumgebung}\label{appendix1:Einrichtung_der_Simulationsumgebung}
\textcolor{red}{TODO:\\
Wie kann man die Simulationsumgebung starten und benutzen und grundlegend einrichten\\
WidowX 200 SDK (modifiziert)
}
